\documentclass{article}

\usepackage{amsmath, amssymb, amsthm}
\usepackage[a4paper, margin=1in]{geometry}

\title{PRIMES in P : The Question}
\author{sid}
\date{\today}

\newtheorem{theorem}{Theorem}
\newtheorem{lemma}{Lemma}
\newtheorem{example}{Example}

\newcommand{\Z}{\mathbb{Z}}
\newcommand{\bld}[1]{\textbf{#1}}
\newcommand{\itl}[1]{\textit{#1}}
\newcommand{\N}{\mathbb{N}}
\newcommand{\Q}{\mathbb{Q}}
\newcommand{\R}{\mathbb{R}}
\newcommand{\eps}{\varepsilon}
\newcommand{\poly}{\mathrm{poly}}
\newcommand{\polylog}{\mathrm{polylog}}
\newcommand{\bigO}{\mathcal{O}}

\begin{document}

\maketitle

\section{The Problem of PRIMES}

Let $n \in \mathbb{Z}$ with $n > 1$. Determine whether $n$ is \textbf{prime}.

\section{Input Model}

For an integer $n$.
\[
\text{Length of n in binary} = k = \lfloor \log_2 n \rfloor + 1
\]

$\implies$ the complexity of algorithms will be measured as a function of $k$, not $n$.

\section{Trial Division is a Lie}

The naivest approach would be to just divide $n$ by every integer from $1$ to $n$.
This should take $O(\sqrt{n})$ time.

However,
$$
O(\sqrt{n})=O(n^{1/2})=O({2^{(k-1)}}^{1/2})=O(2^{\sqrt{k}})
$$

This is \textbf{EXP} time.

Thus, the key question becomes:

\begin{center}
\textit{Is there a deterministic algorithm that decides primality in time polynomial in $k$?}
\end{center}
i.e.,
\begin{center}
\textit{Is PRIMES in P?}
\\
\end{center}

\hrulefill

\subsection*{PS: Wheel Factorization via Primorials}
A simple improvement:
\begin{itemize}
  \item test divisibility by 2
  \item the odd numbers between 3 and $\sqrt{n}$
\end{itemize}

$$
c\#\ \text{(c\ primorial)} = product\ of\ all\ primes \le c
$$

Generalizing further, all primes greater than c\# are of the form $c\#\cdot{k} + i$ for positive integers $i$ and $k$, $0\le{i}<{c\#}$ AND $i$ is coprime to $c$.

\vspace{0.8em}

Consider $6\#=2\cdot3\cdot5=30$
$\implies{\forall{x\ge30}}, x = 30k+i$

Now,
2 divides 0,2,4,6,\dots,28,
3 divides 0,3,9,12,\dots,27,
and 5 divides 0,5,10,15,\dots,25.

Thus, all primes $\ge$ 30 are of the form 30k+i for $i \in{\{1,7,11,13,17,19,23,29\}}$ (but not the other way around, think of 437 = 19 $\cdot$ 23 = 7\#$\cdot$2+17).

$\therefore$ Search space reduced by 76.67\%!!

\vspace{0.8em}

As $c$ grows, the fraction of coprime remainders to remainders decreases, and so the time to test $n$ decreases (though it still necessary to check for divisibility by all primes that are less than $c$).

Observations analogous to the preceding can be applied recursively, giving the Sieve of Eratosthenes.

\hrulefill

\end{document}